\section{Operator overloading}

\begin{frame}{Operator overloading}
    \begin{block}{Operator overloading}
        Operator overloading refers to a single operator's capacity to perform several operations based on the class (type) of operands
    \end{block}
\end{frame}

\begin{frame}{Operator overloading}
    Now in human terms:
    \begin{itemize}
        \pause
        \item Operators like addition (+), equals (==), and greater than (>) can perform certain actions based on their type/class.
        \pause
        \item Other operator are:
        \begin{itemize}
            \item +, -, /, //, \%, ==, in, >, <, >=, <=, !=, and more! 
        \end{itemize}
    \end{itemize}
\end{frame}

\begin{frame}{Operator overloading}
\begin{itemize}
    \item Let's look at some examples
    \pause
    \item Addition:
    \begin{itemize}
        \pause
        \item 5 + 3 = 8 (int)
        \pause
        \item "Cogn" + "AC" = "CognAC" (str)
    \end{itemize}
    \pause
    \item Greater than:
    \begin{itemize}
        \pause
        \item 5 > 3 = True (int)
        \pause
        \item Person(...) > Person(...) = ???
    \end{itemize}
\end{itemize}
\end{frame}

\begin{frame}{Operator overloading}
\begin{itemize}
    \item We can define how these operators should behalve in our classes!
    \pause
    \item We use so called double underscore methods (dunder methods, you may forget that name).
\end{itemize}
\end{frame}

\begin{frame}{What we want}
\begin{itemize}
    \item Consider two objects of our person class: person1 and person2.
    \pause
    \item We want person1 > person2 = True, iff person1 is older than person2.
    \pause
    \item Let's \textbf{overload} the \textit{greater than} operator in our class Person.
\end{itemize}
\end{frame}

\begin{frame}[fragile]{Let's overload}
\begin{minted}[fontsize=\footnotesize]{python}
class Person():
    ...
\end{minted}
\pause
\begin{minted}[fontsize=\footnotesize]{python}
    def __gt__(self, other):  # Note the parameters
\end{minted}
\pause
\begin{minted}[fontsize=\footnotesize]{python}
        return self.age > other.age
\end{minted}
\end{frame}

\begin{frame}[fragile]{Testing never ends}
Let's test again
\pause
\begin{minted}[fontsize=\footnotesize]{python}
person1 = Person("Daan", "Wichmann", 19)
person2 = Person("Damai", "Jiworo", 18)
\end{minted}
\pause
\begin{minted}[fontsize=\footnotesize]{python}
print(person1 > person2)  # What will this return?
\end{minted}
\pause
\begin{minted}[fontsize=\footnotesize]{python}
>>> True
\end{minted}
\end{frame}

\begin{frame}{Why is 'overloading'}
    Why is it called, overloading and not overriding?
    \begin{itemize}
        \pause
        \item We are not overriding operators, because they should still behave the same for integers, strings, etc.
        \pause
        \item We are 'loading' the operator  with functionality for more classes, i.e. \textbf{overloading operators}.
    \end{itemize}
\end{frame}